\pdfoutput=1
\documentclass[11pt]{amsart}
\usepackage{amsmath,amssymb,amsthm,stmaryrd,mathrsfs}
\usepackage[margin=1.2in]{geometry}
\usepackage{xcolor,xurl}
\usepackage[colorlinks=true,linkcolor=blue,citecolor=blue,urlcolor=blue]{hyperref}
\hypersetup{pdfauthor={Bernd Johannes Wuebben},
  pdftitle={A geometric product formula for universal curve counting on Calabi-Yau threefolds}}
\setlength{\emergencystretch}{2em}
\newtheorem{theorem}{Theorem}[section]
\newtheorem{proposition}[theorem]{Proposition}
\newtheorem{lemma}[theorem]{Lemma}
\newtheorem{corollary}[theorem]{Corollary}
\theoremstyle{remark}
\newtheorem{remark}[theorem]{Remark}
\newcommand{\ZZ}{\mathbb Z}
\newcommand{\QQ}{\mathbb Q}
\newcommand{\CC}{\mathbb C}
\newcommand{\PP}{\mathbb P}
\newcommand{\Z}{\mathcal Z}
\newcommand{\cO}{\mathcal O}
\DeclareMathOperator{\Tot}{Tot}
\title[A geometric product formula in universal curve counting]{A geometric product formula for universal curve counting on Calabi--Yau threefolds}
\author{Bernd Johannes Wuebben}
\subjclass[2020]{14N35 (primary); 14J32, 53D45 (secondary)}
\keywords{Universal curve counting, Calabi--Yau threefolds, local curves,
pseudoholomorphic cycles, elementary clusters, virtual fundamental classes}
\date{November 23, 2026\\This version: November 23, 2026}
\begin{document}
\begin{abstract}
We prove a product formula for the universal series of effective
one-cycles on a smooth projective threefold $X$ with trivial canonical
bundle and $H^1(X,\mathcal O_X)=H^2(X,\mathcal O_X)=0$.
The identity holds in Pardon's rational cycle algebra. Its genus-$g$
factor is defined by the circle with opposite weights on two line
bundles of degree $g-1$ whose tensor product is the canonical bundle
of the curve. The exponents are the elementary-cluster integers of
Ionel and Parker, with finite genus support in each integral curve class.
We construct these integers as signed multiplicities in a decomposition
of the cycles on $X$. In the algebra formed from almost-complex
families, a deformation on the sum of two copies of a spin line
identifies each antidiagonal class with the point class of the
corresponding multiple of an elementary curve. Excision for
disconnected cycles factors the image of the series of $X$ at each
energy bound. Stable maps of fixed genus distinguish the multiplicity-one
classes, which suffices to lift the product to the complex cycle algebra.
\end{abstract}

\maketitle
\section{Introduction}\label{sec:intro}

Let $X$ be a smooth projective complex threefold with $K_X\simeq\cO_X$
and $H^1(X,\cO_X)=H^2(X,\cO_X)=0$. For each effective class $\beta$, the characteristic function of the
compact space of cycles representing $\beta$ defines a class $U_{X,\beta}$ in
the Chern-number-zero, virtual-dimension-zero part $H$ of Pardon's rational
algebra of semi-Fano one-cycles \cite[\S\S2--3]{Pardon}. Write
$\Gamma_X\subset H_2(X;\ZZ)$ for the monoid of effective integral classes
and put $U_{X,0}=1$. We use the completion by degree against an ample
integral divisor. A bounded degree contains only finitely many effective
classes, by compactness of the Chow spaces and local constancy of the
integral cycle class.

There are two geometric descriptions of the local curves which enter
this problem. On a smooth curve $C$ of genus $g$, take a line bundle $L$
of degree $g-1$ and set $N=L\oplus(K_C\otimes L^{-1})$. The circle of
weights $+1,-1$ on these two summands fixes the multiple zero sections
$m[C]$. The associated equivariant Thom classes define elements
$a_{g,m}$ of $H$. Write $A_g(q)=1+\sum_{m\ge1}a_{g,m}q^m$.
The second description, due to Ionel and Parker, uses an almost-complex
structure for which every compact curve in a neighborhood of $C$ is
supported on $C$ and the normal operator is injective on every cover.
The contribution of such an elementary cluster is known numerically
\cite[\S3]{IP}. We compare the two descriptions on the space of cycles.

\begin{theorem}\label{thm:main}
For $X$ as above there are unique integers $n_{\beta,g}(X)$ indexed by
$\beta\in\Gamma_X\setminus\{0\}$ and $g\ge0$, zero for all
sufficiently large $g$ at each fixed $\beta$, such that
\begin{equation}\label{eq:product}
 \sum_{\beta\in\Gamma_X}U_{X,\beta}Q^\beta
 =\prod_{\beta\in\Gamma_X\setminus\{0\}}\prod_{g\ge0}
       A_g(Q^\beta)^{n_{\beta,g}(X)}
 \qquad\text{in }H\llbracket\Gamma_X\rrbracket .
\end{equation}
The integer $n_{\beta,g}(X)$ is the signed multiplicity of genus-$g$
elementary clusters of core class $\beta$ in a cluster decomposition
at any energy bound containing $\beta$. It equals the coefficient of
the genus-$g$ elementary series at $\beta$ in the connected
Gromov--Witten expansion of Ionel and Parker.
\end{theorem}

The balanced factors and their operations were constructed in
\cite[\S\S2--3]{PaperIII}. Theorem~1.1 of that paper obtained
\eqref{eq:product} assuming a universal Gromov--Witten evaluation on
analytic families. Here the integer exponents are determined by
deformation of cycles, before taking the connected numerical invariant.
The proof still uses compactness and gluing for stable maps of fixed genus.
Integral homology is used throughout, including when multiplication by
an integer identifies distinct classes.

The local comparison starts with a spin line $L$ and the real rotations
of the two copies of $L$ in $L\oplus L$. These rotations preserve the
Lee deformation of Ionel and Parker. On the integrable fiber the
circle has opposite weights. On each positive fiber the only cycle of
degree $m$ is $m[C]$. A compactly supported class on the total cycle
space restricts to the antidiagonal Thom class at one end and to the
normalized point class at the other (\S\ref{sec:local}).
One can see the effect of the choice of circle already at the double of
a genus-two curve. If $x_{g,m}$ denotes Pardon's scalar class, then
\[
                      a_{2,2}-x_{2,2}=-2x_{4,1}.
\]
In \S\ref{sec:double} we obtain the two negative terms separately from
the two families of branched double covers of a generic balanced curve.

To pass to $X$, we use the cluster isotopy of Ionel--Parker and
Doan--Ionel--Walpuski \cite[\S7]{IP}, \cite[\S3]{DIW}. In a bounded
energy range the disjoint connected supports can be separated in $X$.
A cycle with components in several of the resulting open sets gives a
product of compactly supported classes. At a fold two factors cancel
up to factors of higher multiples of the core class. Iteration gives a
finite signed product in an algebra of almost-complex families
(\S\S\ref{sec:families}--\ref{sec:clusters}).

The difficulty in returning to the complex cycle algebra is that a
relation among classes of curves can run over parameter spaces of
arbitrary dimension. Only the multiplicity-one classes need to be
distinguished. At the least genus in a putative relation the image-cycle
map from stable maps singles out one embedded curve, regular with
local degree one. Sections~\ref{sec:genus} and~\ref{sec:degree} show
that its point orientation survives in compactly supported cohomology.
A comparison of the first coefficient where two products differ then
gives \eqref{eq:product}. The result concerns the balanced antidiagonal
classes. The dependence of the other geometric local elements in
\cite[Remark~6.8]{Pardon} on their defining geometry is a separate question.

\subsection*{Acknowledgements}
I am grateful to John Pardon for a helpful discussion of the local curve classes.

\section{Cycle classes in families}\label{sec:families}

Let $Y\to B$ be a family of almost-complex six-manifolds. The base $B$
is a finite simplicial complex. On each closed simplex the family is
a smooth submersion with a smooth vertical almost-complex structure,
in charts extending across the faces. The restrictions agree on common
faces. Open and empty targets are allowed. Write $Z(Y/B)$ for the space of
compact effective pseudoholomorphic one-cycles in the fibers, including
the empty cycle. Its topology records the base parameter, the Hausdorff
limit of the supports and weak convergence of their integration currents
\cite[Definition~2.6]{DIW}. Empty support is a separate open and closed
locus. A morphism of families is an inclusion of simplicial bases and
a fiberwise almost-complex open embedding into the restricted target,
smooth on every closed simplex.

For a locally compact Hausdorff space $\pi:Z\to B$ put
\[
             V(Z/B)=H_c^0(Z;\pi^*\omega_B),
\]
where $\omega_B$ is the rational dualizing complex. Cohomology here
and below is sheaf cohomology. In cellular form the summand over
an oriented $p$-simplex $\sigma$ consists of the compactly supported
cochains of $Z|_\sigma$ shifted by $p$. The second differential takes
signed restrictions to the faces. This follows by resolving $\omega_B$
by the constant sheaves on the closed simplices with the indicated shifts
\cite[Definitions~2.2 and~3.1, Lemma~3.4]{Pardon}.

\subsection{Topology and operations}
Over a small base chart we can identify a neighborhood of any compact
support with an open part of a product. Parallel transport on one
closed simplex gives such a trivialization after shrinking about the
compact set. A cutoff Hermitian form at the central parameter is
uniformly positive on nearby complex lines, so convergence of currents
bounds their masses in the chart. Consequently changes of trivialization
preserve convergence. Requiring the supports to lie in a relatively
compact neighborhood and imposing a closed bound on mass gives a
compact containing set by \cite[Theorem~2.11]{DIW}. For this application
one may extend the structures by cutoff to the fixed central fiber
on a slightly larger neighborhood. Strict inequalities give a
neighborhood of the prescribed cycle. The finite gluing along closed
face spaces gives a locally compact Hausdorff space over $B$.

On a complex analytic family the topology on each simplex agrees with
Barlet's topology on compact cycles. Apply the support and volume
criterion of \cite[\S1]{Barlet} in the complexification of the
simplex. Only the vertical $(1,1)$ part of a test two-form pairs with
a complex curve. Differences of positive Hermitian forms suffice for
these tests on a compact set. A positive form supported near a given
point of the limiting curve detects lower continuity of supports.
Conversely, compact cutoffs and uniform approximation on a set of
bounded mass give Barlet's continuous volume tests.

Let $S$ be the rational vector-space colimit of $V(Z(Y/B)/B)$ over
these families, identifying subdivisions. A closed base inclusion
includes the corresponding face summands. An open embedding of targets
extends by zero on the open subset of cycles contained in the source.
These two operations commute by restriction of extension by zero.
The external product, with disjoint union of the two targets over the
product base, defines a commutative multiplication on $S$. The empty
target over a point gives the unit. These are the smooth counterparts
of the operations of \cite[\S3.1]{Pardon}. Extension by zero from the
open semi-Fano cycle locus, followed by passage to the underlying
almost-complex family, thus gives a ring homomorphism
\[
                          \Phi:H\longrightarrow S .
\]
On the summand defining $H$ the relative cohomological degree is zero.

If the base is a point, the characteristic function of a compact open
subset of its cycle space gives a class of $S$. Such a class remembers
the ambient almost-complex manifold. An insertion on a family over
$B\times[0,1]$, with $B$ a closed oriented manifold of dimension $n$,
can also be represented by a compactly supported $n$-cocycle.
Restrict it to each top simplex and sum with the product orientations.
In the shifted complex this is a cochain of degree $-1$. Interior
faces cancel in its differential, leaving the difference of the two
endpoint insertions. We use this relation both for characteristic
functions over an interval and for Thom classes over a finite Borel
construction.

For the addition operation, retain all relative degrees in $V$ and
write $S^*$ for the resulting graded group. Proper pullback by addition
over the same base followed by pushforward along the diagonal of bases
takes values in $S^*\otimes S^*$, as in
\cite[Definitions~2.3 and~3.9]{Pardon}. Here the external product is
an isomorphism for the bivariant groups over two independent bases.
Indeed, for the compactly supported cohomology of two locally compact
Hausdorff spaces use their one-point compactifications relative to
the added points. On the compact product finite rectangular covers
are cofinal. The nerve of a product cover maps to the product of the
two nerves by taking marginal barycentric coordinates, with a section
given by their products. Linear interpolation gives a homotopy on
each product of simplices, also on the relative subspaces at infinity.
Finite-polyhedron Künneth over $\QQ$ therefore passes to relative
Čech cohomology, hence to compactly supported sheaf cohomology.
Filtering the two finite face complexes by dimension gives the
bivariant assertion, since the product is an isomorphism on the first
pages. Use product subdivisions containing the diagonal for the
base pushforward. Additivity of support of effective cycles makes
this construction commute with open extension, and restriction to
faces commutes with both operations. The same argument for
multiplication by $d$ gives $\rho_d$.

\section{The class of an elementary curve}\label{sec:local}

For a smooth genus-$g$ curve $C$ and $L\in\operatorname{Pic}^{g-1}(C)$
write $L_+=L$ and $L_-=K_C L^{-1}$. On a cycle of degree $m$ in
$\Tot(L_+\oplus L_-)$, evaluate the elementary symmetric functions of
the two fiber projections at $1+\max(0,m(g-1))$ distinct points of $C$.
The $j$th coefficients are holomorphic sections of $L_\pm^j$, including
at ramification by taking the characteristic polynomial on the finite direct
image from the normalization of each irreducible component and raising
it to the multiplicity of that component. Their vanishing at the chosen points implies
their identical vanishing for $1\le j\le m$. Thus the common zero set is
$m[C]$. In chosen frames the evaluations form an equivariant tuple $F$
of nonzero weights $r_1,\ldots,r_N$. On the associated bundle over
$\PP^N$ it defines a section of the pullback of $\bigoplus_j\cO(r_j)$.
The class
\[
 u_N=(\textstyle\prod_jr_j)^{-1}F^*\operatorname{Th}
                     \bigl(\textstyle\bigoplus_j\cO(r_j)\bigr)
\]
has compact support on the fixed section. Here $\operatorname{Th}$ is
the relative Thom class of the zero section. If $f$ is an additional test and
$i:\PP^N\hookrightarrow\PP^{N+1}$ is a hyperplane inclusion, then
$i_!u_N(F)=h\,u_{N+1}(F)=u_{N+1}(F,f)$, where $h=c_1(\cO(1))$.
The first equality is the hyperplane projection formula and the second
is multiplication by the normalized Euler class of the appended
line. Concatenation of two tuples therefore identifies their universal images
\cite[\S2]{PaperIII}. This construction defines $a_{g,m}$. It applies
also to continuous equivariant tests on an invariant neighborhood when
their common zero set is the same section.

\subsection{Deformation of the spin model}
Fix $L^2=K_C$ and a disk bundle $U$ about the zero section in
$L\oplus L$. Give the two summands the same Hermitian metric. On $U$
use a Kähler form invariant under real rotations of the copies of $L$.
In a spin frame $e^2=dz$ the holomorphic two-form in the construction of
\cite[Proposition~3.2]{IP} is
\[
                  \alpha=dz\wedge(w_1\,dw_1+w_2\,dw_2).
\]
It is the sum of the pullbacks of the tautological two-forms on the
two line summands. The metric and $\alpha$ are rotation invariant.
Hence the Lee almost-complex structures $J_{t\alpha}$ are invariant
for $0\le t\le T$ with $T>0$ small. On the two eigenlines
$w_1\pm iw_2$ the rotations have weights $+1,-1$.

\begin{proposition}\label{prop:local}
For $m>0$, the image of $a_{g,m}$ in $S$ equals the characteristic
function of the single cycle $m[C]$ in the elementary germ.
\end{proposition}
\begin{proof}
Let $Z_m$ be the integrable degree-$m$ cycle space in $U$, and let
$W_m$ be the corresponding space over the Lee interval. Each nonconstant
compact component has positive degree over $C$ by tameness and the
retraction of $U$ onto $C$. Image localization
\cite[Proposition~3.2]{IP} places every such component on $C$ for
$t>0$. Thus $W_m$ is the union of $\{0\}\times Z_m$ and
$[0,T]\times\{m[C]\}$, glued at $(0,m[C])$. Both subsets are closed
in the topology of parameter, support and current. Collapsing the
interval gives a continuous proper equivariant retraction onto $Z_m$.
Its inverse image of a compact set is that set together with the
interval if the set contains $m[C]$.

Pass to the associated bundles for the Hopf circle bundle over $\PP^N$.
At $t=0$ the unit-sphere reduction identifies this with the holomorphic
Borel family restricted to the disk. The pullback of $u_N$ under the
proper retraction has compact support on $\PP^N\times[0,T]$. At $t>0$
it is the normalized Euler class
$(\prod_jr_j)^{-1}\prod_j(r_jh)=h^N$, where $h=c_1(\cO(1))$.
Its integral is $+1$. By the interval relation of \S\ref{sec:families}
the two endpoints give the same class. At the positive endpoint the
Gysin map from a point of $\PP^N$ identifies $h^N$ with the unit at
$m[C]$ in the corresponding target fiber. Excision at the fixed section
permits any smaller disk. The retraction commutes with hyperplane
inclusion, so stabilization respects the entire interval.
\end{proof}

The class in $H$ itself is independent of the balanced splitting.
In a relative Picard chart choose the marked evaluation points as
disjoint local sections and trivialize the evaluation lines. The
fiberwise degree argument gives the same compact section of zeros
over a closed path segment. Its relative Thom class identifies the
endpoints in the analytic diagram under the interval relation.
The variety $\operatorname{Pic}^{g-1}(C)$ is connected. To vary $C$,
take successive charts in connected Teichmüller space, with the
relative degree-$(g-1)$ bundle $\cO((g-1)p)$ for a local marked section
$p$. The fixed-curve comparison identifies choices at common endpoints.
Use piecewise real-analytic paths in the charts, so each segment has
an analytic thickening. For genus zero use the constant sphere and for
genus one use marked elliptic curves. This also joins different spin
choices through balanced splittings.

\subsection{General elementary neighborhoods}
In an elementary neighborhood in the sense of
\cite[Definition~3.1]{IP}, the core is the only irreducible curve,
and its normal operator is the double of a real Cauchy--Riemann
operator $D_0$ on a complex line. Every branched-cover pullback of
this operator is injective. These conditions suffice to identify the
point insertion of the given germ with the class in
Proposition~\ref{prop:local}.

Let $C$ be the core. Choose a tubular chart of its normal bundle
$N$ adapted to the complex splitting along $C$, and pull back
the almost-complex structure $J$ to that chart. On a smaller fixed
disk use $J_s=\delta_s^*J$, $0<s\leq1$, where $\delta_s$ scales
the normal fibers by the positive real number $s$. The family
extends smoothly to the homogeneous almost-complex structure
$J_{\mathfrak d}$ determined by the normal Cauchy--Riemann
operator $\mathfrak d$. Indeed the only block in the pullback with
a negative power of $s$ is $s^{-1}B(x,sv)$ from tangent to normal. Since
$B(x,0)=0$, its extension at zero is obtained by integrating
$\partial_v B(x,\tau sv)[v]$ over $0\leq\tau\leq1$. Identification
of the limit with $J_{\mathfrak d}$ is the normal rescaling formula.
For $s>0$ the elementary isolation condition gives only the cycles
$m[C]$ in this disk. At zero an irreducible curve off the core
would project to a branched cover $f$ of $C$ with a nonzero section
in $\ker(f^*\mathfrak d)$, contrary to super-rigidity. Thus for
each $m>0$ the singleton insertion continues across this ordinary
interval~\cite[Lemmas~3.6 and~3.10 and Proposition~3.12]{DW} and
\cite[Definition 3.1]{IP}.

By the elementary splitting, $\mathfrak d=D_0\oplus D_0$ on
$L\oplus L$, with $D_0$ a real Cauchy--Riemann operator on a
complex line of degree $g-1$. Identify its underlying smooth
complex line bundle with that of a holomorphic spin line. On
this bundle join $D_0$ to the spin Dolbeault operator by affine
interpolation of the zero-order terms. The corresponding normal
models $J_{D_t\oplus D_t}$ project holomorphically to $C$, induce
the usual structure on each fiber and commute with real fiber
dilations and real rotations of the copies. These properties follow
from the construction of $J_{\mathfrak d}$ from the operator. A lift over a holomorphic $f$ is a pair of sections of
$f^*L$ satisfying $f^*D_t$. All models agree as almost-complex
structures along $C$. Thus one fixed Kähler form $\omega$ for the spin model
tames the whole path on a sufficiently small closed disk, by
compactness of $C\times[0,1]$.

Fix $m>0$ and let $r(Z)$ be the maximum fiber norm on the support
of a cycle, using a rotation-invariant metric. For a radius $a>0$
in the interior of the disk the locus $T_m$ of all $r(Z)=a$,
including the parameter $t$, is compact. Its supports lie in a
fixed compact set, and the taming form has integral precisely
$m\int_C\omega$ on each such cycle. Geometric compactness and
continuity of the radius give the claim. There are no circle-fixed
cycles in $T_m$. Each nonconstant irreducible component projects
to a finite holomorphic cover, since a fiber contains no compact
curve. An invariant cycle would therefore have a finite invariant
set in every fiber. The connected rotation group fixes every
element of such a set, and the origin is the only fixed vector~\cite[Theorem 2.11]{DIW}.

On a compact Hausdorff circle space without fixed points there
are finitely many continuous functions of nonzero weights having
no common zero. At a point $x$ its stabilizer is $\mu_k$ for some
$k\geq1$. The function $gx\mapsto g^k$ on its closed orbit
extends continuously to the space by Tietze extension applied to
the real and imaginary parts. Its weight-$k$ Fourier projection
agrees with the original function on the orbit. The nonvanishing
loci so obtained cover the space, hence a finite set suffices.
Apply this to $T_m$. For a test $f$ there, define a test on the
whole degree-$m$ cycle family by

\[
 \widehat f(Z)=
 \begin{cases}
 r(Z)\,f\bigl(t,(\delta_{a/r(Z)})_*Z\bigr),&r(Z)>0,\\
 0,&r(Z)=0 .
 \end{cases}
\]

Homogeneity ensures that the rescaled cycle lies in $T_m$.
Away from zero the formula is continuous by geometric convergence.
At zero it is continuous because $f$ is bounded. The tests retain
their weights and vanish simultaneously precisely at $m[C]$ in
each fiber. If $T_m$ is empty the family already consists of the
singleton in each fiber, and the empty tuple suffices.

Take Borel stage $N$ equal to the number of tests and pull back
their normalized relative Thom class. It has compact support on
the distinguished section over $\mathbb P^N\times[0,1]$. At the
super-rigid end its restriction is $h^N$ and gives the unit at
$m[C]$. At the integrable spin end it gives $a_{g,m}$ because the tests
are continuous circle-equivariant functions, and concatenation
with the holomorphic evaluation tuple after hyperplane stabilization
identifies their normalized Thom classes by the hyperplane formula.

The ordinary rescaling and this Borel deformation give the asserted
identity for the original germ. Write $b_{g,m}=\Phi(a_{g,m})$ and
$b_{g,0}=1$. On the point representatives, pullback by addition counts
the ordered pairs $(r[C],s[C])$ with $r+s=m$. Pullback by multiplication
by $d$ selects $(m/d)[C]$ when $d\mid m$ and is zero otherwise.
Both pullbacks preserve compact supports on these families. For addition,
the support and mass of either effective summand are bounded by those
of the sum, so compactness follows from \cite[Theorem~2.11]{DIW}.
The same reasoning applies to multiplication by $d$.
On these representatives the diagonal of the point base is a point.
The summands are external products of degree-zero classes, giving
\[
  \Delta b_{g,m}=\sum_{r+s=m}b_{g,r}\otimes b_{g,s},
  \qquad
  \rho_d b_{g,m}=
       \begin{cases} b_{g,m/d}&d\mid m,\\ 0&d\nmid m .\end{cases}
\]

\section{Cluster decomposition at bounded energy}\label{sec:clusters}

Choose a Kähler form $\omega$ in an integral ample class and fix a positive
integer $\Lambda$. All series in this section are truncated above
$\omega$-degree $\Lambda$. Initially their labels are zero and all integral
homology classes of positive $\omega$-degree. Write $\mathcal J_{\rm isol}$
and $\mathcal J_*$ for the isolated and regular structures of
\cite[Definition~3.13]{DIW}, respectively. At a structure in
$\mathcal J_{\rm isol}$ the connected supports are embedded irreducible
curves, disjoint from each other. At an energy bound, a cluster consists
of a neighborhood $\mathcal U$ of such a support in the Hausdorff space
of compact subsets, a structure $J$ and an embedded core $C$. The
boundary of $\mathcal U$ misses the bounded connected supports, their
classes inside are positive multiples of $[C]$, and $C$ is the unique
curve of that class there \cite[Definition~3.9]{DIW}. A connected
support carries no multiplicity. The cycle insertions below retain it.

\subsection{Excision for disconnected cycles}
Fix $J\in\mathcal J_{\rm isol}$ and let $\mathscr S$ denote the space of nonempty connected supports at $J$
of degree at most $\Lambda$. It is compact by
\cite[Theorem~2.11 and Lemma~3.16]{DIW}. Indeed connectedness is closed
under Hausdorff convergence and every nonempty cycle has integral energy
at least one. Fix a finite partition of $\mathscr S$ into relatively open
and closed subsets $\mathscr S_i$. The union $K_i$ of all curves in
$\mathscr S_i$ is compact by the closed incidence relation in
$\mathscr S_i\times X$. The sets $K_i$ are pairwise disjoint since distinct
irreducible curves are disjoint. Choose open neighborhoods $V_i$ of $K_i$
with disjoint closures. For every sufficiently close tame $J'$ each
bounded connected support is contained in some $V_i$. Otherwise a
sequence of such supports violating the assertion would have a connected
limit at $J$ inside some $K_i$, contrary to Hausdorff convergence. For
a prescribed finite partition by cluster neighborhoods with unoccupied
boundaries, the same argument shows that the partition by the $V_i$
continues to agree with it near $J$. Include the complement of the
selected cluster in the initial partition.

For a union $\mathscr U$ of parts define $P(\mathscr U,J)$ to have
constant term one. Its coefficient at $\beta\ne0$ is the characteristic
function of the cycles in class $\beta$ all of whose connected
components are supported in the corresponding open neighborhoods. Use
the class in $H_2(X;\ZZ)$ under inclusion for cycles in an open target.
These locally constant functions have compact support. To see this over
a short closed path, use compactness at fixed energy and the continuity
of the integral homology label \cite[Theorem~2.11 and Remark~2.13]{DIW}.
Once the $V_i$ are separated, restricting the integration current to
one of them amounts to multiplying by a cutoff which is one on its
closure and zero near the other closures. The nonempty supports converge
there in the Hausdorff topology. Occupation of the several neighborhoods
is locally constant. Thus the insertion for a union of parts extends
over the path, and the interval relation preserves its class in $S$.
These short transports may use smooth segments in charts of tame
structures. Compactness and avoidance of the boundaries hold in a
neighborhood of the given structure, so segments in a smaller convex
coordinate neighborhood join the same endpoints even if the original
isotopy path has only $C^1$ parameter dependence.

At either endpoint, open excision into $\coprod_i V_i$ decomposes each
cycle uniquely as $\sum_i D_i$, including $D_i=0$. In total class
$\beta$ the characteristic function is the sum of external products
over all $\sum_i\beta_i=\beta$. There are finitely many relevant
labels by compactness. Consequently
\begin{equation}\label{eq:excision}
                P(\mathscr U,J)=\prod_{i:\,\mathscr S_i\subset\mathscr U}
                                    P(\mathscr S_i,J).
\end{equation}
Shrinking the $V_i$ about the $K_i$ gives the same insertions by open
excision. The separation in the target is what identifies the right
side with multiplication in $S$.

\subsection{Crossings of the wall}
For a core in class $A$ consider the subgroup of the units of the
truncated series algebra generated by the cluster factors with core
class $dA$, $d\ge2$, at isolated structures. Equivalence modulo this
subgroup is an exact identity after multiplication by finitely many
such factors and their inverses. Every factor has constant term one.

The subgroup is generated by corrections with unobstructed cores.
Indeed, for a cluster around a possibly obstructed core in class $dA$,
keep its outer neighborhood and its complement fixed as above. Choose
a nearby structure in $\mathcal J_*$ in the neighborhood given by
\cite[Propositions~3.14 and~3.21]{DIW}. Transport its insertion along
a short path. At the new endpoint the outer neighborhood partitions
into finitely many clusters with unobstructed cores of classes $kdA$,
$k\ge1$, by Proposition~3.21(2) of that work. Include the complementary
part and apply \eqref{eq:excision} to this partition. This expresses
the entire correction as the product of the resulting regular factors,
even if the former core disappears.

Consider first a transverse fold of class-$A$ cores
\cite[Proposition~3.33]{DIW}. The whole path lies in
$\mathcal J_{\rm isol}$, including the fold parameter, where the
obstructed core is isolated. Spatial separation there extends the
outer cluster insertion across the fold. At the two-core endpoint isolate
the two degree-one branches and partition the remainder. At the other
endpoint there are no class-$A$ cores. Equation~\eqref{eq:excision} at
both ends gives
\[
             P_+P_-\prod_i H_i=\prod_j K_j,
\]
where each $H_i,K_j$ has core class a proper multiple of $A$. Thus the
two distinguished factors are mutual inverses modulo those corrections.
This product accounts also for cycles meeting both core neighborhoods.

Along a regular arc of cores in the isolated locus, the factor is
constant modulo the same subgroup. Use the finite cover by outer
neighborhoods in \cite[Proposition~3.30]{DIW}. On an overlap choose
a common smaller cluster around the core
\cite[Proposition~3.19]{DIW}. Each difference is a union of clusters
in proper multiples, and \eqref{eq:excision} compares the factors.
At a wall along a path fixing the core, use the two-parameter model
$s=x(x-t)$ of \cite[Proposition~3.37]{DIW}. Along $s=0$ the two branches
are $x=0$ and $x=t$. At small $s>0$ the root
$(t+\sqrt{t^2+4s})/2$ joins the distinguished branch at $t=-\delta$
to the other branch at $t=\delta$. For $s<0$ the two positive roots
at the latter end meet at the transverse fold $t=2\sqrt{-s}$.
Use isolated auxiliary paths as in that proposition. The endpoint
insertions persist under the small displacements in $s$ by spatial
separation at the two original endpoints. Regular transport on the
first arc and the fold relation on the second therefore invert the
distinguished factor at the fixed-core crossing.

\subsection{The elementary endpoint}
For an embedded core $C$ with induced complex structure $j$, choose
a spin line $L$ and give $L\oplus L$ the Kähler form
$\Omega=\pi^*(\omega|_C)+i\partial\bar\partial\|w\|^2$ on a small
disk, using equal Hermitian metrics. The symplectic normal bundle
of $C$ has first Chern number $2g-2$. Classification of complex
rank-two bundles over a surface gives a unitary normal isomorphism
with $L\oplus L$. Choose it as the derivative of a tubular embedding
using the symplectic orthogonal to $TC$. The forms agree as bilinear
forms along $C$. Linear interpolation on a smaller disk and the
relative Moser argument, with primitive obtained by fiber contraction,
identify the symplectic forms there. Extend the compatible integrable
structure from a smaller tube by extending its metric and using polar
decomposition with $\omega$. Inside the tube introduce the Lee
deformation with small amplitude $T\chi$, where $\chi$ is one on an
inner disk $U$ and supported in the region of agreement. The resulting
ambient $J_L$ is tame and agrees with the elementary structure on $U$
\cite[Proposition~3.2 and proof of Theorem~7.1]{IP}.

Take $C\subset V\Subset U$ a smaller disk bundle. At $J_L$ no cycle
contained in $\overline V$ meets $\partial V$, since every irreducible
component there lies on $C$. Compactness at bounded energy excludes
such boundary supports for all sufficiently close structures.
Containment in $V$ thus gives a compact open and closed insertion over
any sufficiently short closed path. At $J_L$ it equals the point
insertion of Proposition~\ref{prop:local}. At a nearby generic
structure fixing the unobstructed core, as in
\cite[Proposition~3.54]{DIW}, it is the cluster insertion in a small
Hausdorff ball about $C$. Indeed a sufficiently small ball implies
containment in $V$, and otherwise a sequence of connected supports
in $\overline V$ outside that ball would limit to $C$. This proves
the comparison inside the original compact target.

\begin{proposition}\label{prop:finite}
Through degree $\Lambda$ the image of the compact series in $S$ is a
finite product of elementary factors $\Phi(A_g(Q^\gamma))$ and their
inverses. Each occurrence carries the genus and integral class of its core.
\end{proposition}
\begin{proof}
Start with a generic structure and the finite partition of
\cite[Proposition~3.18]{DIW}. Follow the generic fixed-core path and its
finite subdivision in the proof of \cite[Theorem~3.25]{DIW} to the
elementary endpoint just identified. Regular arcs transport the factor,
and each crossing inverts it modulo the proper-multiple factors.
The sign is $(-1)^w$, where $w$ counts the fixed-core crossings relative to the elementary endpoint. Keep the corrections in the exact identities and replace
each by its regular factors as above. For each recursive correction
the positive integer $\lfloor\Lambda/\langle[\omega],A\rangle\rfloor$
strictly decreases. The branching and the depth are finite. Finally
the insertion for the \emph{entire} $\beta$-locus extends from the
generic structure to the given projective structure over a tame path
by compactness at bounded energy and continuity of the integral class.
The endpoint relation in $S$ completes the proof.
\end{proof}

\section{Detection by maps of fixed genus}\label{sec:genus}
Write $x_{h,m}$ for the scalar classes of
\cite[Definition~4.7]{Pardon}. The map $\Phi$ of \S\ref{sec:families}
will only have to distinguish scalar classes of multiplicity one.

\begin{theorem}\label{thm:independence}
The images $\Phi(x_{h,1})$, $h\ge0$, are linearly independent over $\QQ$.
\end{theorem}

We first reduce an arbitrary linear relation to a statement about a
single moduli space. Section~\ref{sec:degree} proves the required
nonvanishing there.

\subsection{The image cycle}
Let $Y\to B$ be a family as in \S\ref{sec:families}. Let $M_k$ be
the coarse moduli space of connected unmarked stable maps in the fibers,
of arithmetic genus $k$, with nonconstant image and vertical Chern number
zero. Write $z_k:M_k\to Z(Y/B)$ for the map taking the image cycle.
A component is counted with its degree over the normalization of its
image. Contracted components contribute zero.
\begin{lemma}\label{lem:proper}
The map $z_k$ is continuous and proper.
\end{lemma}
\begin{proof}
Gromov convergence gives convergence of the weighted currents and the
Hausdorff supports. Contracted components have image in the nonconstant part of the
image because the domain is connected. For properness, take a
sequence above a compact subset of $Z(Y/B)$ and pass to a subsequence whose
cycles converge. Work over one closed simplex, in a local target
trivialization about the limiting support. Hausdorff convergence puts
all the images in one relatively compact neighborhood. The masses of
the cycles bound the energies of the maps, with multiplicity included.
Indeed one can test against a cutoff Hermitian form of the limiting
parameter, uniformly positive on all nearby complex lines in this
neighborhood. For the fixed parametrizing surface in
Ivashkovich--Shevchishin's theorem take a closed oriented surface
of genus $k$. Each connected nodal domain of arithmetic genus $k$
is the quotient of such a surface by pinching disjoint circles,
as one sees by smoothing its nodes. The boundary condition in
their theorem is empty. The images stay in a common compact set,
and the structures on its neighborhood converge smoothly after
trivialization. Their theorem now gives a convergent subsequence
of stable maps. On smooth parts away from the nodes elliptic
regularity upgrades convergence to the smooth Gromov topology for
these smoothly converging structures. If the limit map were constant,
all nearby images would lie in a coordinate ball on which a fixed
exact two-form tames the nearby structures. Its integral over a
nonconstant closed curve is positive, whereas Stokes' theorem gives
zero. Thus the limit is nonconstant and its image cycle is the
prescribed limit. The empty cycle has a neighborhood disjoint from
all nonempty cycles, and the integer Chern number is locally constant
in the Gromov topology. These observations preserve the restrictions
defining $M_k$. The compactness argument applies in metrizable local
charts for the two topologies and over a finite union of simplices.
For compactness see \cite[Theorem~1.1]{IS}, and for the image-cycle map see \cite[Definition~2.14]{DIW}.
\end{proof}

\subsection{A regular map over the least genus}
Use a smooth curve $C$ of genus $h$ and a degree $h-1$ line $L$
outside the theta divisor. For $h=0$ take $\mathcal O(-1)$ and for
$h=1$ a nontrivial line of degree zero. Serre duality and
Riemann--Roch give $H^0(C,N)=H^1(C,N)=0$ for
$N=L\oplus K_C L^{-1}$. Every compact degree-one cycle in $N$ is a
section. Indeed it has one irreducible component mapping with degree
one to $C$, no compact curve lies in an affine fiber, and the
normalization gives the graph of a holomorphic section. Hence the
core is isolated in that degree and the point insertion represents
$x_{h,1}$, as in \S\ref{sec:local}.

For a stable map over this insertion, exactly one normalized
component maps with degree one to $C$ and it maps isomorphically.
Arithmetic genus equals the sum of the component genera plus the
first Betti number of the dual graph. It is at least $h$. In the
equality case the dual graph is a tree and every other vertex is a
contracted rational component. Any nontrivial such tree has a
contracted leaf, violating stability in the absence of markings.
Thus $z_h^{-1}([C])$ is the single embedded map, with trivial map
automorphisms in every genus. Its unparametrized deformation and
obstruction spaces are $H^0(C,N)$ and $H^1(C,N)$, since the
domain tangent complex cancels the tangential part of the target
complex. The map is regular with complex orientation $+1$.

In this splitting the whole degree-one cycle locus is the singleton
$[C]$. For both circles its unit represents the local class, by the
empty equivariant tuple at Borel stage zero. In particular
$a_{h,1}=x_{h,1}$. For $k<h$ the fiber of $z_k$ over this cycle is empty.

\begin{lemma}\label{lem:relation}
Suppose a finite relation among $\Phi(x_{h,1})$ has least genus $k$
with nonzero coefficient. There is a smooth family of almost-complex
six-manifolds over an open subset of $\CC^r$ for which the point
orientation at one regular map, in a product neighborhood of a
marked parameter, has zero image in $H_c^{2r}(M_k;\QQ)$.
\end{lemma}
\begin{proof}
Represent the sum on a finite disjoint
union of vertex targets, one per nonzero coefficient. In this
comparison diagram the filtered-quotient argument of Pardon's Lemma
3.2 applies. Take disjoint unions for common targets. For parallel
inclusions, attach a cylinder on the source base between their two
images. Use the corresponding pulled-back larger target near each
end and the source open target across the middle. Stacking three
cylinders permits simplicial attachment even when the endpoint
images overlap. The resulting target is Hausdorff since near the middle slice only the common open target is present, with
larger pieces on opposite sides. It is simplexwise smooth.
The two inclusions induce equal maps after attachment by the
constant-cylinder prism on compactly supported cochains and open
extension. A subdivision likewise has a triangulated product
concordance. Thus every finite vanishing relation is witnessed in
one family $X/B$ after these attachments. Distinct marked vertices
of the initial disjoint-union base remain distinct under its
inclusion~\cite[Lemmas 3.2--3.4]{Pardon}.

At each marked vertex $b$ attach a new closed interval to $B$ by
one endpoint only, and use the constant family with fiber $X_b$
on that interval. Its other endpoint $b'$ has an open neighborhood
in the enlarged base on which the family is a product.
The original base is a subcomplex, so its zero relation still
vanishes after inclusion. The compact point insertion crosses the
new edge by the interval relation for a constant family. Indeed
the insertion extended constantly over the edge has the difference
of its vertex insertions as its face boundary. Use distinct new
endpoints for the different marks. Relabel the enlarged complex
$B$ and use these endpoints as the marked parameters below.

Embed a subdivision of $B$ as a subcomplex in an oriented Euclidean
space and take a compact PL regular neighborhood, with $B$ in its
interior $W$. After subdivision there is a PL neighborhood retraction
onto $B$ that fixes $B$. Pull the family back along this retraction. The inclusion
of $B$ takes the zero relation to zero by the bivariant pushforward.
Equivalently, the dualizing complex of the neighborhood restricts
in the interior to $\mathbb Q[n]$, and the closed-base pushforward
has support on $B$ in that interior. One can use the pushforward
directly with target $W$, by properness of the closed inclusion,
so no restriction of compact supports across an arbitrary open set
is needed. Increase the ambient dimension if necessary. On $W$
the relation is therefore zero in $H_c^n(Z(X_W/W);\mathbb Q)$.
This is the base thickening used in the proof of Pardon's Theorem
5.15. It requires only a PL retraction for the present families~\cite[proof of Theorem 5.15]{Pardon}.

The closed-base pushforward for $i:B\hookrightarrow W$
uses the counit $i_*\omega_B=i_*i^!\omega_W\to\omega_W$.
Pull back this counit to the cycle space and use proper base change
for the closed inclusion of its restriction over $B$. This formula
also gives functoriality for a vertex followed by $i$, and commutes
with pulling back to the stable-map space by the proper map $z_k$.

Continuity of the retraction gives a Euclidean ball about each new
endpoint mapping into its constant-family neighborhood. Over this
ball the targets, cycle spaces and stable-map spaces are products.
The singleton cycle is open in the original open target and remains
open in the larger fiber by open embedding of the cycle spaces.
It is compact there as well. The oriented point Gysin map in the
ball is the external product of this singleton insertion and the
orientation at the marked parameter. Pullback by $z_k$ computes
the two factors separately. At a higher-genus mark the inverse image
of the singleton is empty by the genus bound above. At the least-genus mark it consists
of the one regular embedded map, so over the ball it is precisely
a copy of the ball. The pulled-back class is its point orientation
extended by zero, with a nonzero coefficient. This proves the reduction.

It remains to smooth the dependence of the targets on the parameter.
Let $P$ denote the compact triangulated PL neighborhood just constructed,
with interior $W$ open in Euclidean space. Increase its dimension to
an even integer $2r$. The family on $P$ is simplexwise smooth after
subdivision. Using the attached constant intervals, choose
disjoint Euclidean balls $B_{4\delta}(b)\Subset W$ around the
marked parameters on each of which the entire family is a product.
Use a common $\delta>0$ for these finitely many marks. The point
orientation on the distinguished stable-map graph maps to zero in
$H_c^{2r}(M_k)$ for a putative dependence. Subdivide to mesh
smaller than $\delta/4$.

For every vertex $v$ of the subdivision take its open star in $P$.
Choose a partition of unity $\lambda_v$ with closed supports contained
in these stars, smooth on the interior in the Euclidean smooth
structure. It can be obtained by restricting a smooth partition of
unity from an open Euclidean cover near $P$. Define $F:P\to P$ by
the barycentric combination of the vertices with weights
$\lambda_v(y)$. At $y$, the vertices whose closed supports contain
$y$ belong to the least simplex containing $y$. Since the cover is
finite, all other weights vanish in a whole neighborhood of $y$.
Consequently $F$ maps that neighborhood into a \emph{single} simplex
and its barycentric coordinates there are smooth in $W$. The
straight-line homotopy from $y$ to $F(y)$ stays in its simplex and
moves points by at most the mesh. In particular it preserves the
boundary of $P$.

Modify $F$ to fix the marked balls. Choose a smooth function
$\chi$ equal to zero on the balls of radius $\delta$ and one
outside the balls of radius $2\delta$. Replace $F(y)$ by
$(1-\chi(y))y+\chi(y)F(y)$, choosing $0\leq\chi\leq1$. This
map remains in the least simplex containing $y$, moves points by
less than $\delta/4$, and agrees with the old map near the boundary
of $P$. Where it changes, its source and image lie in one of the
balls of radius $3\delta$ with constant target family. Outside
the modifying balls the earlier factorization through one simplex
applies. Write $F$ for the modified map.

Set $W'=F^{-1}(W)$. The restriction $W'\to W$ is proper, as a base
change of the continuous map of compact spaces $P\to P$, and $W'$
lies in the Euclidean interior of $P$. Pull the target family back
along this map. It is a smooth submersion with smooth vertical
almost-complex structure on $W'$. On the modifying balls it is a
product. Elsewhere it is the pullback of a smooth simplex model.
On overlaps of simplex models the maps factor through the common
face and the prescribed identifications are smooth. The product
trivializations are the specified constant families on the attached
intervals, so they agree with the simplex models there. For the stable-map
spaces this is a Cartesian pullback, including their Gromov topologies,
so their induced map is also proper.

For a marked $b$, the preimage of $B_{\delta/2}(b)$ equals that
ball and $F$ is its identity. Indeed any preimage lies within
$3\delta/4$ of $b$ by the displacement bound. Consequently the
proper pullback of the point orientation on the distinguished graph
is the very same point orientation extended by zero. It still
vanishes in the global compactly supported group.
\end{proof}

\section{The local virtual degree}\label{sec:degree}

The putative relation of Lemma~\ref{lem:relation} would kill a
regular point orientation in the presence of other, possibly singular
and noncompact, components of the moduli space. We use an implicit atlas
on a neighborhood of the compact support of such a vanishing. First
we construct its pairing with compactly supported cohomology. We then
obtain the charts, including their orientation at nodal maps, from
fixed-target gluing. In this section references to numbered axioms,
deformation spaces and gluing charts are to \cite{Atlas}.

\subsection{A functional on a finite union of regular footprints}
In an implicit atlas, $X_I$ denotes the thickening belonging to a
finite set $I$ of indices, with obstruction space
$E_I=\bigoplus_{\alpha\in I}E_\alpha$, section $s_I:X_I\to E_I$,
and group $\Gamma_I=\prod_{\alpha\in I}\Gamma_\alpha$.
The footprint is the image of the zero locus under the quotient map
to the moduli space. The regular footprint uses the regular part of
the thickening. The local orientation system is the virtual
orientation system of the atlas.

\begin{proposition}\label{prop:functional}
Let $M$ be locally compact Hausdorff. Suppose we have the chart,
footprint and transversality data of an oriented topological implicit
atlas of virtual dimension $n$ on a neighborhood of a compact set.
We use no boundary on these charts. Let $O$ be an open set covered
by regular footprints of a finite set of indices $A$, allowing
multi-indices formed from $A$. All coefficients are rational.
There is a functional

\[
                   H_c^n(O)\longrightarrow\mathbb Q
\]

for this choice of charts. Suppose a point $p\in O$ has an oriented
manifold neighborhood of dimension $n$ in the zero locus. For some
index set $I\subset A$ assume its zero-coordinate lifts near $p$
are regular, are freely permuted in fibers of size $|\Gamma_I|$,
and that $s_I$ on this regular thickening is a submersion there.
Orient the neighborhood of $p$ by the virtual orientation. The
functional sends the point orientation class extended by zero to
$+1$.
\end{proposition}
\begin{proof}
For $I\subset J\subset A$ take the deformation
space $Y_{IJ}=X_{I,J,A}$ of Pardon's Definition 4.2.1. It is a
manifold with corners of dimension $n+\dim E_A+|I|$. Write
$e:Y_{IJ}\to E_A$ for its obstruction projection. On $e^{-1}(0)$
there is a continuous footprint map to the unthickened space.
Delete the part of $e^{-1}(0)$ lying over the complement of $O$,
an ordinary closed subset, and call the resulting open manifold
$Y^O_{IJ}$. Write $Z_{IJ}$ for its full zero locus and
$\pi:Z_{IJ}\to O$ for the footprint map. These restrictions commute
with the maps (4.2.1) for the one fixed set $A$. Away from $e=0$
no map to $O$ is required.

Let $O^+$ be the one-point compactification (with a disjoint point
if $O$ is compact). To a compact $K\subset O^+$ associate the
complex of relative singular chains

\[
 F_{IJ}^\bullet(K)=
 C_{n+\dim E_A-\bullet}
 \big(Y^O_{IJ},Y^O_{IJ}\setminus\pi^{-1}(K\cap O);
                     o_{E_A}^{\vee}\big)^{\Gamma_A}.
\]

Here $\pi^{-1}(K\cap O)$ is closed in $Y^O_{IJ}$, including when
$K$ contains infinity. Use the normalized transition maps of
Definition 4.2.6 and take their finite homotopy colimit $F^\bullet$.
This is the one-point compactification argument of Pardon's proof
of Lemma A.6.4 applied to each of the deformation spaces.

We give the descent checks. The assignments on compact subsets
are homotopy K-sheaves. The value at the empty set is zero, and the
two-compact-set axiom is singular-chain excision for the open
complements. For decreasing compact neighborhoods, a singular chain
has compact image. Its intersection with $Z_{IJ}$ is compact, and
the image of that intersection in $O^+$ is compact. If the chain
misses $\pi^{-1}(K\cap O)$, its footprint therefore misses a compact
neighborhood of $K$ as well. This proves continuity under shrinking
neighborhoods, including at infinity. Finite group invariants are
exact over $\mathbb Q$. Finite homotopy colimits preserve these
conditions.

At each $x\in O$, the stalk calculation is the one in Lemma 4.3.1
and Proposition 4.3.3. Excision away from the deleted zero loci
does not change the relative chains at $x$. The footprint diagrams
give $o_{M,x}$ in degree zero and zero in every other degree.
Coverage is required only over $O$. At infinity every relative
chain complex is zero. For the uniform lower bound in the purity
criterion, use Poincaré--Lefschetz duality with closed support in
the finite-dimensional manifolds $Y^O_{IJ}$, as in Lemma A.6.4.
The relative homology vanishes above their dimensions. There are
finitely many pairs and finitely many nondegenerate strings in
their indexing poset. Thus their shifted homotopy colimit still
has a uniform lower bound, also on neighborhoods of infinity.
The resulting degree-zero sheaf is $j_!o_O$ for $j:O\hookrightarrow
O^+$, by adjunction and the stalk calculations. Proposition A.5.4
now identifies $H^*(F^\bullet(O^+))$ with $H_c^*(O;o_O)$~\cite[\S\S4.2--4.3, Proposition A.5.4 and Lemma A.6.4]{Atlas}.

At $K=O^+$ the complement in each pair is the complement of the
\emph{whole} zero set in $Y^O_{IJ}$. Projection to $E_A$ therefore
defines a map to the relative chains of $(E_A,E_A\setminus0)$.
In the convention where the transition $J\to J'$ is the ordinary
pushforward divided by $|\Gamma_{J\setminus J'}|$, take the
projection on the $IJ$ term divided by $|\Gamma_J|$. Indeed
$|\Gamma_J|=|\Gamma_{J'}||\Gamma_{J\setminus J'}|$, and the
obstruction projection is unchanged by the transition. This gives
a chain map from the homotopy colimit by the augmentation on
strings of length zero, with value zero on higher strings. In
degree $n$, identify the resulting obstruction-space orientation
class with $1$. This gives the functional. Applying the projection
directly to chains relative to a proper compact subset of a larger
zero set would lack this map of pairs.

\par\smallskip\noindent\emph{Normalization at the regular point.} Let $V_I\subset O$ be the regular
footprint for $I$, intersected with $O$. In the diagonal term
$I=J$, take the slice $t_i=1$ for $i\in I$. It is the open part
of $E_{A\setminus I}\times X_I^{\rm reg}$ obtained by deleting
zeros over the complement of $O$. Its obstruction projection is
$(e',x)\mapsto(e',s_I(x))$. Its entire zero set is closed in
this open manifold and has proper finite quotient onto $V_I$.
Indeed it is the restriction of the finite group quotient of the
zero-coordinate chart to its invariant open regular locus.

Here is the slice calculation directly from Definition 4.2.1.
Denote the slice by $D$ and its zero set by $Z_D$. Before removal
of the zeros outside $O$ the closed set over a compact $K\subset O$
in this slice is

\[
       \{0\}\times\big((s_I|_{X_I^{\rm reg}})^{-1}(0)
                       \cap\psi_{\emptyset I}^{-1}(K)\big)
       \ \subset\ E_{A\setminus I}\times X_I^{\rm reg}.
\]

The pair uses the complement of this entire set. This detail
matters in using (4.3.14), whose displayed complement takes the
product of $E_{A\setminus I}$ with the complement in the second
factor. For example, take a point with trivial groups,
$X_\emptyset=\{p\}$ regular, and one additional chart
$X_\alpha=\mathbb R$ with $s_\alpha=\mathrm{id}$. For
$I=J=\emptyset$, Definition 4.2.5 gives
$H_1(\mathbb R,\mathbb R\setminus\{0\})=\mathbb Q$.
The pair $(\mathbb R,\varnothing)$ would give zero in degree one.

For any compact $K\subset O^+$ put $L=\pi^{-1}(K\cap O)$ in
$Y=Y_{II}^O$. On the interior $Y^\circ$, all $t_i>0$ and
$(t,e',x)$ are product coordinates, with $e_i=s_i(x)/t_i$.
In these coordinates the pair $(Y^\circ,Y^\circ\setminus L)$
is $(\mathbb R_{>0})^{I}$ times
$(D,D\setminus(L\cap D))$. Retraction of the first factor to
$(1,\ldots,1)$ proves the slice isomorphism on the interior.
The inclusion of this interior pair into $(Y,Y\setminus L)$
induces an isomorphism as well. To see it, apply Poincaré--Lefschetz
duality (A.6.6) to the closed set $L$ in $Y$ and to its intersection
with $Y^\circ$ in the open manifold $Y^\circ$. The coefficient
sheaf on $L$ in the first duality is the extension by zero of the
orientation sheaf from $L\cap Y^\circ$. Inclusion of chains is the
corresponding extension on compactly supported cohomology. This
naturality follows from the relative-chain map in the proof of
Lemma A.6.4, which is the orientation identification on interior
stalks and has zero source at boundary stalks. Both dual groups are
therefore the same. Tensoring with $o_{E_A}^\vee$ and taking
finite-group invariants preserves the isomorphism. The argument
also covers $K=O^+$, where $L$ is the whole zero set.

Applying Lemma A.6.4 once more to $Z_D$ in the slice, then taking
finite-group invariants as in (4.3.7), identifies the diagonal
term on $O^+$ with $H_c^*(V_I;o_O)$. The map of this term into
the homotopy colimit induces the identity of orientation sheaves
on $V_I$ by Proposition 4.3.3. Its source sheaf is extension by
zero from $V_I$. By adjunction this specifies the map of sheaves
on all of $O^+$, so on cohomology it is precisely extension by
zero $H_c^*(V_I;o_O)\to H_c^*(O;o_O)$.

Take a ball at $p$ on which the zero-locus cover splits into
$|\Gamma_I|$ sheets. Pullback of its point orientation to the
cover gives coefficient one on each sheet. In the relative chains
of the slice these classes are small oriented normal disks at the
lifts. Submersivity lets us choose them with boundary disjoint
from the entire zero set and obstruction projection of local degree
one. In chains with obstruction orientation coefficients, first
average one disk over $\Gamma_{A\setminus I}$, which fixes its
lift and acts simultaneously on the complementary obstruction
coordinates and their orientation coefficients. Its projected
degree remains one. Sum its translates over $\Gamma_I$ to obtain
the invariant representative on the free sheets. The
previous paragraph identifies its image in the full complex with
the point orientation extended by zero. Projection of the orbit sum
and division by $|\Gamma_I|$ therefore give one. This also works when regularity of
the unperturbed \emph{fixed-domain} operator fails. It uses an ordinary
submersion in one nonempty chart over the smooth ball.
\end{proof}
If the local orientation at $p$ vanished in $H_c^n(M)$, it would
already vanish with support in a compact set $K_0$ containing
$p$. Choose finitely many regular footprints covering $K_0$ and
an open $O$ inside their union containing $K_0$. Excision for
support in $K_0$ would make its image in $H_c^n(O)$ zero. The
functional contradicts this.

\subsection{An implicit atlas for maps in the fibers}
Let $\pi:E\to W$ be the smooth family of Lemma~\ref{lem:relation},
with $W$ open in $\CC^r$ and $n=2r$. Choose a smooth horizontal distribution
and give $E$ the almost-complex structure given by the vertical
structure and the lift of the standard structure on $W$.
Every connected compact holomorphic map to $E$ projects to a
constant. Holomorphic functions on each normalized component are
constant and agree at nodes. Thus unperturbed fixed-target maps
are exactly the maps into the fibers, including their Gromov
topology. Their Chern number in $E$ equals their vertical Chern
number since $TW$ is trivial.

The fixed-target construction can be used separately in each
homology class of maps to $E$. This class is locally constant in
the labeled-domain topology: on nearby domains collapse the
plumbing necks to the nodes, compare maps into a small tubular
neighborhood of the limiting graph, and use short geodesics in
$E$ to homotope to the limiting map composed with the collapse.
The collapse sends the fundamental cycle to the sum of the
fundamental cycles of the limiting components. Equivalently one
can integrate the corresponding integral currents in a compact
neighborhood and use local constancy of the integral current
homology class in the flat topology. A compact set of maps meets
only finitely many such sectors. The test also restricts to
nonconstant maps. A constant map has an image neighborhood
in a coordinate ball with an exact taming form, which excludes a closed
nonconstant holomorphic curve in that ball.

Use the indices of \cite[\S9.2]{Atlas} for the fixed target $E$ with \emph{vertical}
perturbation sections. Use actual codimension-two incidence disks
in $E$ and the usual marked-domain orbifold charts. The positivity
of rank of the differential on every unstable nonconstant component
gives the stabilizing disks of Lemma 9.2.8. In a small target
trivialization one can thicken a transverse disk in the fiber in
all base directions and round its edge away from the marked
intersections. For each finite index set $I$, take the fixed-target
thickening with these vertical perturbations. Every solution, at
arbitrary vertical obstruction coordinates, still projects
holomorphically to $W$ and lies in a single fiber. The spaces, group
actions, footprint quotients, obstruction coordinates and forgetful
maps are therefore the ordinary data of Definitions 9.2.1--9.2.4.
The finite group need not act effectively on its marked-domain chart.

Choose indices on sufficiently small marked-domain charts. On each
one we can also specify a complex $\Gamma_\alpha$-representation
$P_\alpha$ of dimension $kr$, where $k$ is the fixed arithmetic
genus, and equivariant perturbation forms whose projection induces
an isomorphism

\[
           P_\alpha \ \simeq\
             H^1(T,\mathcal O_T)\otimes\mathbb C^r
\]

on each labeled domain $T$ in the chart. Here is the construction.
At the central stabilized curve use the right-hand space as the
representation. Smooth $(0,1)$ forms supported in the smooth locus
away from the marked points span its first factor by duality with
dualizing differentials. Such a differential vanishing on the smooth
locus is zero, including at every branch of a node. Choose a complex
linear right inverse and average over the finite group. The forms
extend on disks in the smooth locus of the universal curve. Extend
and average again, keeping the original equivariant values at the
center. The classes still give a basis after shrinking the chart.
Indeed the dual Hodge bundle is locally free for the proper flat
family of connected nodal curves of fixed genus. In a Hodge frame
the pairing with these disk forms is continuous up to the central
fiber, by integration over a fixed compact set away from the nodes.
The extension may be chosen smooth on the total smooth locus.
This also handles normalization components on which the map is
constant. Take the horizontal lifts of the resulting
$\mathbb C^r$-valued forms on $\mathcal C_\alpha\times E$. They
are allowed perturbations in Definition~9.2.1(vi) of \cite{Atlas}. For relative duality and base change see \cite[Tags 0E6N and 0E6S]{Stacks}.

For example, take the rational nodal genus-one curve obtained by
identifying $0$ and $\infty$, with a stabilizing mark at $1$. The
involution $\tau(z)=1/z$ acts by $-1$ on the dualizing differential
$dz/z$, hence by $-1$ on $H^1(\mathcal O_T)$. Choose a form
$\phi$ supported on a disk near $2$ pairing nontrivially with
$dz/z$. The form $(\phi-\tau^*\phi)/2$ is anti-invariant and has
the same nonzero pairing. An invariant horizontal coordinate
would miss this cokernel. Taking $P_\alpha$ with its actual
stabilizer representation deals with this node without passing to
an unlabeled quotient.

For each nonempty $I$ select for the moment $\alpha\in I$. Replace
the datum at $\alpha$ by the same datum with obstruction space
$E_\alpha\oplus P_\alpha$ and the additional horizontal perturbations.
Keep every other datum unchanged. This is an ordinary collection
of Pardon's fixed-target indices. On any simultaneous thickening
containing $\alpha$ the projected perturbed equation reads

\[
        \bar\partial w+\eta_\alpha(h)=0,\qquad
        w=\pi u,\quad h\in P_\alpha .
\]

Taking its Dolbeault class gives $h=0$, and then $w$ is constant.
Conversely every solution of the original vertical thickening
is a solution with $h=0$. This is a homeomorphism of the two
thickenings. All marks and domain labels are identical, the graph
topologies agree, and the extra coordinate on one side is constantly
zero. If the distinguished index is retained, the homeomorphism
also intertwines the footprint maps on the locus where the new
obstruction coordinates are zero.

The regularity test in Definition 9.2.5 is independent of which
$\alpha\in I$ is chosen for this purpose. At a solution the
horizontal block of the fixed-domain linearization is
$(\xi,h)\mapsto\bar\partial\xi+\eta_\alpha(h)$. Its cokernel is
zero, and its kernel is the constant horizontal displacement
$T_bW$. Derivatives of the horizontal forms with respect to domain
labels have coefficient $h=0$. In the kernel of this block the
remaining equation is exactly the vertical linearization with
free base value $b$. Transverse intersections with incidence disks,
used to define the labels in (9.2.3), move smoothly with the map
on the fixed domain. Thus the full augmented operator is surjective
if and only if the vertical operator with free base is surjective.
Triviality of the automorphism group is the same on both sides.

Define regularity on nonempty thickenings by this intrinsic test,
and set the regular locus of the empty chart to the empty set.
For $I\subset J$ with $I$ nonempty, augment \emph{one index in $I$}
on both sides of the inclusion. Under the two homeomorphisms just
proved, fixed-target Proposition 9.2.6 gives the required openness
and the topological submersion onto $E_{J\setminus I}$. The
auxiliary space never occurs in this quotient. For $I$ empty the
submersion and regular-footprint freeness requirements have empty
domains of verification. The footprint and quotient axioms still
hold there without a regularity assumption. The dimension in the
augmented model is

\[
       2r(1-k)+\dim_{\mathbb R} E_I+2rk=n+\dim_{\mathbb R}E_I
\]

for unmarked maps in the Chern-zero sector.

For coverage choose a stabilizing disk and finite-dimensional
vertical perturbations spanning the fixed-domain \emph{vertical}
cokernel, as in the proof of Lemma 9.2.9. At the unperturbed map
the label derivatives of these vertical perturbations have zero
coefficients. Perturbations supported away from all special points
suffice by that lemma's elliptic argument. Extend the sections
locally in the vertical tangent bundle and use the induced finite
group representation for equivariance. The choice of a labeled
domain chart removes the residual automorphisms of the lifted object.
This produces a nonempty regular footprint at every point. The construction
uses compactness of the domain curves, and only finitely many indices
are needed near a compact $K_0$. Pardon's compact-target assumption
in the global GW application supplies compactness of the entire
unperturbed space. Proposition 9.2.6 itself is a gluing calculation
near a single compact map. Its stabilized fixed-target model in
Appendix B.1 uses a smooth almost-complex target and transverse
disks. All the maps in this local calculation stay near the fixed compact image. Local compactness of the unperturbed
space also follows from the regular footprints themselves. Each is
an open set in the Hausdorff moduli space and the finite-group
quotient of the closed zero locus in a locally compact regular
chart. These footprints cover.

\subsection{Orientation on overlaps} At a smooth-domain regular
point of $X_I$, retain the same domain labels and incidence disks.
For $\alpha,\beta\in I$, write $H_T=H^1(T,\mathcal O_T)\otimes
\mathbb C^r$ and denote the two maps to this quotient by
$A_\alpha:P_\alpha\to H_T$ and $A_\beta:P_\beta\to H_T$.
Identify $P_\alpha$ with $P_\beta$ by the complex isomorphism
$G=A_\beta^{-1}A_\alpha$. On the common linearized domain consider
the horizontal perturbation forms

\[
            \eta_s=(1-s)\eta_\alpha+s\,\eta_\beta G,
            \qquad 0\leq s\leq1 .
\]

Every $\eta_s$ projects to $A_\alpha$. Thus
$(\xi,h)\mapsto\bar\partial\xi+\eta_s(h)$ is onto for the whole
path, and its kernel always has $h=0$ and $\xi$ constant. This
also describes the horizontal equation with domain variations,
since the projection of the original map is constant and the
horizontal perturbation coefficient at the solution is zero.
On this kernel the remaining linear equation, including the old
perturbations and moving incidence marks, is the same for every
$s$. Regularity therefore gives a path of surjective linearizations
with the very same tangent kernel. The canonical determinant
orientation varies continuously along that path in the fixed
kernel. Identifying its two auxiliary factors by $G$ preserves
their real orientations. This proves equality of the two
orientations of the regular thickening at the smooth point.
Only a path of linear operators at that point is needed. The
construction commutes with relabeling, since $G$ is defined by
the induced maps to the same horizontal cohomology group.

Each augmented chart already has a continuous orientation
identification by Proposition 9.2.6(iii). Both live on the same
topological regular locus by the homeomorphisms above. By \cite[\S B.12]{Atlas}
the plumbing parameters can all be nonzero arbitrarily close to
every nodal point of that locus. The two identifications therefore
agree there as well, since their ratio is a locally constant sign.
For $I\subset J$ with $I$ nonempty, choose an index in $I$ on both
sides. On the regular footprint the forgetful map at zero extra
coordinates is a finite covering of an open set in $X_I^{\rm reg}$.
Smooth-domain lifts are dense on this covering. On these lifts the
ordinary determinant exact sequence for adding $E_{J\setminus I}$
gives the orientation transition for the vertical charts, since
the same oriented $P_\alpha$ cancels from both sides. Continuity
gives the transition at the nodal lifts. For a triple of nested
nonempty index sets one may choose the same index in the smallest
one throughout, so the transition maps compose. All extra complex
factors have even real dimension. This verifies orientation descent
on the regular nodal overlaps used by the compact-support test~\cite[Proposition 9.2.6 and \S\S B.12--B.13]{Atlas}.

\begin{proof}[Proof of Theorem~\ref{thm:independence}]
Use the regular point over the least genus in Lemma~\ref{lem:relation}.
Its normal bundle is acyclic, its map automorphism group is trivial,
and the family is constant over a ball at its marked parameter.
In any nonempty regular chart covering this point the map to the
vertical obstruction coordinates is a submersion near its lift.
Indeed the unperturbed linearization with domain variations is onto.
On this smooth domain the implicit function theorem solves the perturbed
equation with arbitrary prescribed small obstruction coordinates.
Transverse intersections recover the stabilizing marks. The kernel
is the tangent to the graph of the unperturbed solutions, oriented
by the complex normal operator and the base. There are $|\Gamma_I|$
lifts. Proposition~\ref{prop:functional} therefore pairs the local
point orientation to $+1$ on a neighborhood of every compact set
containing it. The vanishing of Lemma~\ref{lem:relation} is impossible.
\end{proof}

\section{Lifting the product to the complex cycle algebra}\label{sec:product}

\subsection{Stabilization of the integers}
At cutoff $\Lambda$, let $N^\Lambda_{\gamma,h}$ be the sum of the signs
of the elementary factors of genus $h$ and core class $\gamma$ in
Proposition~\ref{prop:finite}. These integers are zero outside a finite
set of pairs. Suppose two decompositions, possibly one obtained by
truncating a decomposition at a higher bound, give different integers.
Cancel their common lower-degree factors and take a class $\beta$ in
the first degree of disagreement. Because every nonconstant term of a
factor in class $\gamma$ has degree at least that of $\gamma$, comparison
of the coefficients of $Q^\beta$ gives
\[
           \sum_h (N^\Lambda_{\beta,h}-\widetilde N_{\beta,h})\,
                       \Phi(x_{h,1})=0 .
\]
Here we used $a_{h,1}=x_{h,1}$. Theorem~\ref{thm:independence} gives
equality in every genus, a contradiction. The same induction proves
that $N^\Lambda_{\beta,h}=0$ when $\beta$ is ineffective for the
original projective structure. The compact series has coefficient
zero there, and sums of effective lower labels are effective.

For the chosen polarization denote the stabilized counts by
$n_{\beta,h}$. Taking $\Lambda$ equal to the degree of $\beta$ proves
finite genus support at $\beta$. Thus Proposition~\ref{prop:finite}
has already given the product of Theorem~\ref{thm:main} after
application of $\Phi$. We need only multiplicity-one independence to
lift this identity.

\subsection{Primitive classes}
Let $P=\QQ[X_{h,m}:h\ge0,\ m\ge1]$ and $X_{h,0}=1$. The map
$p:P\to H$ taking $X_{h,m}$ to $x_{h,m}$ is surjective by the
geometric generation theorem \cite[Theorem~6.13]{Pardon}.
The Chern numbers of the generators of virtual dimension zero are
nonnegative, so only generators of Chern number zero can enter
this summand. The scalar classes respect the coproduct $\Delta$ from
addition of cycles and pullback $\rho_d$ under multiplication by $d$.
Their formulas are
\[
 \Delta X_{h,m}=\sum_{r+s=m}X_{h,r}\otimes X_{h,s},
 \qquad
 \rho_d X_{h,m}=
 \begin{cases}X_{h,m/d}& d\mid m,\\ 0& d\nmid m.\end{cases}
\]
For the scalar classes, apply the division formula and the diagonal
calculation of \cite[\S4, in particular Lemma~4.13]{Pardon}. Each term
with a nonzero shift has one tensor factor of negative virtual dimension
and thus vanishes by \cite[Theorem~6.13]{Pardon}. The same vanishing
removes the shifted terms in the balanced calculation
\cite[\S3]{PaperIII}. Thus both operation formulas use only the
geometric generation and dimension statements on the complex side. The augmentation on both algebras is
restriction to the empty cycle.

In the triangular coordinates
$v_{h,m}=[q^m]\log(\sum_{j\ge0}X_{h,j}q^j)$ the generators are
primitive and $\rho_d v_{h,m}=v_{h,m/d}$ for $d\mid m$, zero otherwise.
Suppose the kernel of $p$ contains a nonzero polynomial $f$ of smallest
total polynomial degree in these coordinates. Its constant term is zero
by augmentation. Its reduced coproduct lies in the tensor square of
polynomials of strictly lower degree. The map $p$ is injective on
this vector space and remains so on its tensor square. Since the
reduced coproduct maps to zero, $f$ is primitive. Over $\QQ$ an
additive polynomial is linear, by differentiating $f(v+w)=f(v)+f(w)$
at $w=0$. Applying $\rho_j$ for the largest multiplicity $j$ occurring
in $f$ gives a nonzero formal combination of the $X_{h,1}$ in the
kernel. Their images even under $\Phi p$ are independent by
Theorem~\ref{thm:independence}. Hence $p$ is an isomorphism.

\begin{lemma}\label{lem:primitive-division}
Let $\operatorname{Prim}(H)$ denote the primitive subspace for $\Delta$.
Then
\[
  \operatorname{Prim}(H)\cap\bigcap_{d>1}\ker\rho_d
      =\bigoplus_{h\ge0}\QQ x_{h,1}.
\]
\end{lemma}
\begin{proof}
In the coordinates $p(v_{h,m})$ a primitive is an additive polynomial,
hence linear. If its largest multiplicity is $j>1$, applying $\rho_j$
leaves exactly the coefficients at multiplicity $j$, now on the
linearly independent classes $x_{h,1}$. Conversely, each $x_{h,1}$
is primitive and is killed by every $\rho_d$ with $d>1$.
\end{proof}

\subsection{The first coefficient of disagreement}
Write $F$ for either the compact series or the product with the
stabilized exponents, formed in $H\llbracket\Gamma_X\rrbracket$.
Each has $F_0=1$ and satisfies
\begin{equation}\label{eq:ops}
 \Delta F_\beta=\sum_{\gamma+\delta=\beta}F_\gamma\otimes F_\delta,
 \qquad
 \rho_d F_\beta=\sum_{\substack{\gamma\in\Gamma_X\\d\gamma=\beta}}F_\gamma .
\end{equation}
For the compact series take inverse images of the entire class locus
under addition and multiplication. For the other series use the
formulas on the balanced elements, multiply, and take inverses for
negative exponents. The substitution $Q^\gamma\mapsto Q^{d\gamma}$
is multiplicative even when several exponents have the same image.
Each effective root occurs separately in \eqref{eq:ops}.

If $R_\beta$ is a nonzero difference at a least ample degree, the
terms with two positive summands cancel in its coproduct. It is
primitive. For $d>1$ all the roots on the right of \eqref{eq:ops}
have smaller degree, so $\rho_d R_\beta=0$. By
Lemma~\ref{lem:primitive-division}, $R_\beta$ is a linear combination
of the $x_{h,1}$.
The finite-energy identity says $\Phi(R_\beta)=0$.
Theorem~\ref{thm:independence} now forces $R_\beta=0$.
This proves the product formula. Uniqueness follows by comparing the
first degree of disagreement between any two systems of exponents.
A second polarization yields an identity in the same coefficients. For either polarization there are finitely many effective sublabels
of any given $\beta$, and their possible decompositions are finite.
Uniqueness therefore identifies the two systems.

\subsection{Identification with the elementary-cluster coefficients}
It remains to identify the integers already obtained with those of
the connected Gromov--Witten expansion. Apply fixed-target connected
counts to the \emph{geometric decompositions} of
\S\ref{sec:clusters}. For a compact open insertion $u$ let
$\epsilon(u)$ be its value at the empty cycle, and let $D_k(u)$ be
its localized positive-energy connected genus-$k$ count. The same
open part of the stable-map space computes this count on excision
into a target neighborhood. For the compact symplectic targets
and paths in that section the localized counts are additive on
partitions and invariant under transport
\cite[Lemma~3.8 and Proposition~3.10]{DIW}. On disjoint targets a
connected map lies in just one of the two, hence
\[
                D_k(uv)=D_k(u)\epsilon(v)+\epsilon(u)D_k(v).
\]
Extend the equality coefficientwise with the integral class labels.
Each of our normalized series has empty-cycle value one as a series.
Thus connected counts add under their multiplication and change sign
on inversion, which can equally be calculated in the truncated
series algebra. This applies along the exact fold, subdivision and
endpoint relations used to obtain each finite product.

For genus $h$ the elementary connected contribution is
\[
 G_h(q,t)=\log\!\left(1+
 \sum_{m\ge1}q^m\sum_{\lambda\vdash m}
       \prod_{\square\in\lambda}
           \bigl(2\sin(\ell(\square)t/2)\bigr)^{2h-2}\right),
\]
where $\ell(\square)$ is the hook length \cite[Proposition~3.3]{IP}.
Thus the signed occurrences of each factor give the coefficients of
$G_h(Q^\beta,t)$ in the connected expansion of $X$. Uniqueness is
\cite[Theorem~3.50]{DIW}. At the first degree of disagreement, the
least remaining genus is distinguished by
$[q]G_h(q,t)=t^{2h-2}+O(t^{2h})$. This identifies the integers of
\eqref{eq:product} and completes the proof of Theorem~\ref{thm:main}.

\section{A double curve of genus two}\label{sec:double}

In genus two the first difference between the two circle actions can
be calculated inside the complex cycle algebra. The identity also
follows under numerical reconstruction in \cite[\S6]{PaperIII}.

\begin{proposition}\label{prop:double}
In $H$ one has $a_{2,2}=x_{2,2}-2x_{4,1}$.
\end{proposition}
\begin{proof}
Let $C$ be a smooth genus-two curve with hyperelliptic involution
$\iota$. Choose distinct non-Weierstrass points $p,q$ in distinct
orbits. Choose $L$ with $L^2=\mathcal O(p+q)$, possible by
surjectivity of multiplication by two on the Jacobian. Put
$M=K_C L^{-1}$. The divisors of chosen sections $s_L,s_M$ of
$L^2,M^2$ are $p+q,\iota p+\iota q$, respectively. Indeed
$p+\iota p\sim q+\iota q\sim K_C$. Each section spans its space
of global sections by Riemann--Roch, since neither divisor is
canonical. A section of $L$ would give a square section with a
double zero, which is impossible. The same holds for $M$. Also
$ML^{-1}=K_C(-p-q)$ is a nontrivial degree-zero line bundle.
By the deformation of balanced pairs in \S\ref{sec:local} and the
scalar construction \cite[\S4.3]{Pardon}, both classes can be
computed on this splitting.

In $Y=\operatorname{Tot}_C(L\oplus M)$, every irreducible compact
curve projects finitely onto $C$. A degree-two cycle has zero trace
in both fiber coordinates $u,v$, because $H^0(L)=H^0(M)=0$.
Its two fiber points, with multiplicity, thus have the form
$\pm(u,v)$. The three quadratic invariants give global sections

\[
        u^2=a s_L,\qquad v^2=b s_M,\qquad uv=d\in H^0(K_C).
\]

They extend over ramification by symmetric evaluation on the finite
maps from the normalizations of the components. Since
$d^2=ab\,s_Ls_M$ and the four zeros of $s_Ls_M$ are simple,
$ab=0$ and $d=0$. Conversely each point of the two coordinate
axes gives a cycle. On the first axis it is
$D_a=\{u^2=a s_L,\ v=0\}$, and there is an exchanged formula
on the second. The origins are the same cycle $2[C]$. These
parametrizations are continuous, including at zero, by the finite
flat families of spectral divisors. Their inverse is continuous by
fiber evaluation away from the zeros. This identifies the underlying
cycle space equivariantly with the union of the two complex lines.
For nonzero amplitude each cover is smooth and connected, with
two simple branch points. Riemann--Hurwitz gives genus four. The
two independent scalings of $L,M$ act on $a,b$ with characters
$(2,0),(0,2)$.

\smallskip\noindent\emph{The two Borel classes.}
Set $B=\mathbb P^1\times\mathbb P^1$ and let $h_1,h_2$ be the
positive hyperplane classes. Twist $L,M$ by the hyperplane lines
from the first and second factors, respectively. The relative
degree-two cycle space $W$ is the union of the total spaces
$P_+=\operatorname{Tot}\mathcal O(2,0)$ and
$P_-=\operatorname{Tot}\mathcal O(0,2)$, glued along their zero
sections $B$. Write $U_+,U_-$ for the complex Thom classes of
these two line bundles. Use the two base inclusions

\[
            i_s(z)=(z,z),\qquad i_a(z)=(z,\bar z),
\]

orienting the source by its complex orientation in both cases.
The first restricts the target family to the scalar Borel family.
On the second, the second hyperplane line identifies real analytically
and fiberwise complex linearly with $\mathcal O(-1)$ by the
Hermitian metric. Thus its restriction is the antidiagonal Borel
family. A real analytic base map and real analytic bundle
identification suffice here. On real simplices they complexify to
the family germs of \cite[Definition~3.1]{Pardon}. The Gysin map of either
embedding preserves the universal image. One can use a product
triangulation subdivided along the diagonal, and transport it by
conjugation of the second factor for the second graph. Both give
the canonical concordance class of immersed real analytic
triangulations of $B$. In the shifted compactly supported
description the pushforward is the oriented normal Thom map, as
in Pardon's diagonal computation~\cite[\S3.1, Lemma 3.5 and proof of Lemma 4.13]{Pardon}.

At Borel stage one the scalar cut is $(a+b)^*\mathrm{Th}/2$,
where the sum is in $\mathcal O(2)$. For the antidiagonal cut
use $a+j(b)$, where the Hermitian conjugate-linear identification
$j:\mathcal O(-2)\to\mathcal O(2)$ reverses the real fiber
orientation. Again divide by two. Each cut vanishes exactly on the
common zero section. These choices represent the published
normalizations. To check this even for the continuous test $j$,
use the standard Hermitian metrics at every finite Borel stage.
The tests restrict along hyperplanes, and each ordinary normalized
Euler class is $h$. Append a defining tuple from \cite[\S2]{PaperIII} to the
one-test tuple and append this test to the defining tuple. The two
resulting compactly supported Thom products coincide at the common
stage by naturality and the projection formula of the stabilization
lemma. The same argument identifies the scalar test with that of \cite[Definition~4.7]{Pardon}.

The Poincaré duals of the two base graphs are $h_1+h_2$ and
$-h_1+h_2$, respectively. For the second sign, integrate $h_1$
and $h_2$ on the graph to obtain $1,-1$. Naturality of the
normal Thom map under restriction to either branch gives the
following table in $H_c^4(P_\pm;\mathbb Q)$.
\[
\begin{array}{lcc}
 & P_+ & P_- \\ \hline
\text{scalar} & (h_1+h_2)U_+/2 &(h_1+h_2)U_-/2\\
\text{antidiagonal} & (-h_1+h_2)U_+/2 &(h_1-h_2)U_-/2\\
\text{difference} &-h_1U_+ &-h_2U_-
\end{array}
\]
These restrictions determine the class on $W$. For the proper
map $\nu:P_+\sqcup P_-\to W$ the difference of restrictions
to $B$ gives the exact sequence
$0\to\mathbb Q_W\to\nu_*\mathbb Q_{P_+\sqcup P_-}
\to\mathbb Q_B\to0$. There are no higher direct images since the
fibers are finite. Now $H^3(B)=0$, so restriction of $H_c^4(W)$
to the two branches is injective.

For the positive class $(h_1U_+,0)$, fix one point on the first
base factor, trivialize its hyperplane line there, and fix a nonzero
amplitude $a_*$. Over the remaining base $\mathbb P^1$, take
the complex point class of $a_*$ in the amplitude line and extend
by zero from a disk avoiding zero. Pushing forward from the base
slice gives precisely $(h_1U_+,0)$ by the Thom isomorphism. Its
support misses the other branch. The exchanged construction
represents $(0,h_2U_-)$. It remains to compute the universal
image of the class on each base slice.

\smallskip\noindent\emph{The class of a branch.}
We treat the first branch. The family over the remaining
$B'=\mathbb P^1$ has first summand $L$ and second summand
$M\otimes\mathcal O_{B'}(1)$. Write $S=\operatorname{Tot}_C L$.
Deform the threefolds as affine line bundles over $S$, retaining
the projection to $S$. Choose independent sections
$\sigma_0,\sigma_1$ of $\mathcal O_{B'}(1)$. Let $z_0$ be
the simple zero of $\sigma_0$. Choose a nonzero class

\[
     c\in H^1(C,ML^{-1})\simeq\mathbb C .
\]

Cover $C$ by a disk $V$ avoiding the zeros of $s_L$, and the
complement $U$ of its center. These two Stein sets and their
intersection compute this cohomology. Represent $c$ on $U\cap V$.
For $t\in[0,1]$, add the following translation to the usual
transition of the second fiber coordinate, written intrinsically
in the common line over the overlap. The added term is

\[
     t\,c\,u\,
          \big[\sigma_0+(u^2/s_L-a_*)\sigma_1\big].
\]

This is holomorphic on the whole overlap, including $u=0$.
The two-set gluing with inverse translation defines a holomorphic
family, also for complex $t$. The relative tangent determinant
in each threefold is the pullback of $TC\otimes L\otimes M$, up
to a constant line on the fiber, so the cycles stay semi-Fano.

Fix $a$ in a small disk about $a_*$ avoiding zero. Any cycle
of degree two with this amplitude projects onto the smooth
spectral divisor $D_a\subset S$ with degree one. Indeed its
projected cycle is determined by $u^2=a s_L$, and no component
can lie in an affine line fiber. The finite degree-one map to the
smooth divisor is an isomorphism on the normalization and the lift
is a section of the affine line bundle over $D_a$. Conversely
every such section gives the embedded cycle. Write
$\pi:D_a\to C$. The identity
$\pi_*\mathcal O_{D_a}=\mathcal O_C\oplus L^{-1}$ gives

\[
   H^0(D_a,\pi^*M)=0,\qquad
   H^1(D_a,\pi^*M)=H^1(C,ML^{-1}) .
\]

The lifting obstruction is
$t[\sigma_0(z)+(a-a_*)\sigma_1(z)] [c\,u]$.
Here multiplication by $u$ identifies $c$ with its nonzero
class in the odd summand of the displayed direct image. Hence a
lift exists precisely when this expression vanishes. When it does,
the translations on the divisor are themselves zero, and the zero
section in the two local charts gives the unique lift.

Pull back the Thom class of $a-a_*$ on this amplitude disk. It
defines a supported class throughout the prism $B'\times[0,1]$.
Its zero set consists of the section $D_{a_*}$ over all of $B'$
at $t=0$, together with $D_{a_*}$ over $z_0$ for every $t$.
This set is compact. The amplitude is continuous by fiber
evaluation in the projection to $S$. On the zero-obstruction
locus the explicit local zero graphs vary continuously by their
spectral equations. Thus the parametrization above is also a
homeomorphism over the disk. Extension by zero and the oriented
prism relation identify the endpoint insertions. At $t=1$, the
relative cycle space near the support is the graph

\[
                  a=a_*-\sigma_0(z)/\sigma_1(z).
\]

The function on the right has nonzero derivative at $z_0$.
The insertion has local degree $+1$ on this graph. Projection
of the graph to the base is a complex orientation preserving local
isomorphism, so the Gysin image from the fiber over $z_0$ of the
unit on the isolated curve is exactly this class.

For regularity in that fiber, the normal sequence of $D=D_{a_*}$
has subbundle $\pi^*M$ and quotient $\pi^*L^2$. The latter is
the normal line to $D$ in $S$, by its divisor equation.
Its space of sections is generated by the amplitude variation
$da=1$, since $H^0(L^2)=\mathbb C s_L$ and $H^0(L)=0$.
The normal connecting map sends this variation to
$\pm\sigma_1(z_0)[c\,u]$, by differentiating the transition
on $D_a$ at $a=a_*$. It is an isomorphism. Also
$H^1(D,\pi^*L^2)=H^1(L^2)\oplus H^1(L)=0$. It follows that
$H^0(N_{D/Y_{z_0,1}})=H^1(N_{D/Y_{z_0,1}})=0$.

Finally deform the ambient germ along this smooth curve to the
total space of its normal bundle $N$. One model is the blow-up
of $Y_{z_0,1}\times\mathbb C$ along $D\times0$, with the
strict transform of the central fiber removed. The copy of $D$
persists and is isolated and regular in every fiber, including
the normal bundle fiber. Near a smooth embedded curve of multiplicity
one the relative cycles are embedded deformations, and vanishing
of both normal cohomology groups makes the relative space locally
the base. Thus the section over a closed parameter interval is
compact and open in the relative cycle space there, and its unit
continues to the zero section of $N$. Every degree-one compact
cycle in $N$ is a holomorphic section, so its degree-one cycle
space is the single point $H^0(D,N)=0$. In \cite[Definition~4.7]{Pardon} one can take Borel stage zero and the empty tuple,
since the whole degree-one locus is compact and its class is the unit.
To identify the indices, adjunction gives
$\deg N=2g(D)-2=6$, since the fiber threefold has trivial
canonical bundle. Thus $c_1(T\operatorname{Tot}N)[D]=0$.
Pardon's scalar construction applies to any rank-two holomorphic
bundle on a smooth proper curve of this genus and Chern number.
This includes the possibly nonsplit $N$ obtained from the normal
sequence. With $m=1$ the empty tuple at stage zero therefore gives
$+x_{4,1}$, with the sign fixed by the local degree on the graph~\cite[Remark 4.1 and Definition 4.7]{Pardon}.

Exchanging $L$ and $M$ preserves the assumptions, including
nontriviality of the degree-zero ratio, so the second positive
branch class has the same image. Substitution in the last row of
the table gives $-2x_{4,1}$ for the difference.
\end{proof}

\begin{thebibliography}{Wue26}
\bibitem[Bar99]{Barlet} D.~Barlet,
\emph{How to use the cycle space in complex geometry},
in \emph{Several Complex Variables}, Math. Sci. Res. Inst. Publ.
\textbf{37} (1999), 25--42,
\href{https://library.slmath.org/books/Book37/files/barlet.pdf}{MSRI Publications}.
\bibitem[DIW26]{DIW} A.~Doan, E.-N.~Ionel and T.~Walpuski,
\emph{The Gopakumar--Vafa finiteness conjecture},
Ann. of Math. (2) \textbf{204} (2026), 819--871,
\href{https://doi.org/10.4007/annals.2026.204.3.1}{doi:10.4007/annals.2026.204.3.1}.
Theorem numbering follows \href{https://arxiv.org/abs/2103.08221v2}{arXiv:2103.08221v2}.
\bibitem[DW21]{DW} A.~Doan and T.~Walpuski,
\emph{Castelnuovo's bound and rigidity in almost complex geometry},
Adv. Math. \textbf{379} (2021), 107550,
\href{https://doi.org/10.1016/j.aim.2020.107550}{doi:10.1016/j.aim.2020.107550}.
\bibitem[IP18]{IP} E.-N.~Ionel and T.~H.~Parker,
\emph{The Gopakumar--Vafa formula for symplectic manifolds},
Ann. of Math. (2) \textbf{187} (2018), 1--64,
\href{https://arxiv.org/abs/1306.1516v6}{arXiv:1306.1516v6}.
\bibitem[IS99]{IS} S.~Ivashkovich and V.~Shevchishin,
\emph{Gromov compactness theorem for stable curves}, preprint (1999),
\href{https://arxiv.org/abs/math/9903047}{arXiv:math/\allowbreak 9903047}.
\bibitem[Par16]{Atlas} J.~Pardon,
\emph{An algebraic approach to virtual fundamental cycles on moduli
spaces of pseudo-holomorphic curves}, Geom. Topol. \textbf{20} (2016),
779--1034, \href{https://doi.org/10.2140/gt.2016.20.779}{doi:10.2140/gt.2016.20.779}.
\bibitem[Par25]{Pardon} J.~Pardon,
\emph{Universally counting curves in Calabi--Yau threefolds},
\href{https://arxiv.org/abs/2308.02948v3}{arXiv:2308.02948v3} (2025).
\bibitem[Stacks]{Stacks} The Stacks Project Authors,
\emph{The Stacks Project}, Algebraic Curves,
\href{https://stacks.math.columbia.edu/tag/0E6N}{Tag 0E6N} and
\href{https://stacks.math.columbia.edu/tag/0E6S}{Tag 0E6S}.
\bibitem[Wue26]{PaperIII} B.~J.~Wuebben,
\emph{Antidiagonal local curves and integral exponents in universal curve counting},
second edition, 22 November 2026,
\href{https://doi.org/10.5281/zenodo.24432700}{doi:10.5281/\allowbreak zenodo.24432700}.
\end{thebibliography}

\end{document}
